% GENERATED FILE. Edit canonical lessons, then run npm run generate:books.
% canonical-source-hash: fnv1a64:9574471c60ce2e89
% canonical-lessons: ES-C518-robar, ES-C518-robo, ES-C518-sospechoso, ES-C518-atracar, ES-C518-huir

\chapter{Un robo en el barrio}
\label{ch:es-un-robo-en-el-barrio}
\hlchaptermodality{518}

\begin{chapteropening}
\textbf{By the end of this chapter:} \emph{I can report a robbery: what was stolen, who the suspect is, and how the thief ran away.}

Complete the last lesson of chapter 518: I can report a robbery: what was stolen, who the suspect is, and how the thief ran away.
\end{chapteropening}

\section[robar]{robar — to steal, to rob}

\label{lesson:ES-C518-robar}



\begin{quote}
Before the new one: say the Spanish for to notice, then the Spanish for the detail.
\end{quote}

\subsection*{You'll want to know: robar}
\textbf{robar} — \textquotedblleft{}to steal, to rob\textquotedblright{}.

\begin{quote}
Me robaron la cartera en el metro.
\end{quote}

\begin{quote}
\emph{Someone stole my wallet on the metro.}
\end{quote}

\emph{Robar algo a alguien} — to steal something from someone.

\begin{cousinweb}
\emph{robar} comes from a Germanic verb \emph{*raubōn} 'to rob, to plunder', through Medieval Latin \emph{raubāre} (EtymDB). English \emph{rob} goes back to the same Germanic verb.
\end{cousinweb}

\begin{grammarlens}[title={What you've built}]
One more word for telling, explaining and describing at B1.
\end{grammarlens}

\subsection*{Your turn}
\begin{itemize}
\raggedright
  \item \emph{Say it:} \emph{robar}
  \item \emph{Say it:} \emph{robar}, once more
  \item \emph{Say it:} a sentence of your own with \emph{robar}
  \item \emph{Recall:} say the Spanish for to notice, then the Spanish for the detail, then say \emph{robar} again
\end{itemize}

\begin{tcolorbox}[breakable,colback=teal!4,colframe=teal!35!black,title={Before you move on}]
What does robar mean? (\textquotedblleft{}To steal, to rob\textquotedblright{}.) Say it once more.
\end{tcolorbox}
\section[el robo]{el robo — the robbery, the theft}

\label{lesson:ES-C518-robo}



\begin{quote}
Before the new one: say the Spanish for the detail, then the Spanish for to steal.
\end{quote}

\subsection*{You'll want to know: el robo}
\textbf{el robo} — \textquotedblleft{}the robbery, the theft\textquotedblright{}.

\begin{quote}
Hubo un robo en la tienda de la esquina.
\end{quote}

\begin{quote}
\emph{There was a robbery at the shop on the corner.}
\end{quote}

\emph{¡Esto es un robo!} — this is daylight robbery! — about a high price.

\begin{cousinweb}
\emph{robo} is built on the verb \emph{robar} (EtymDB), which comes from a Germanic verb \emph{*raubōn} 'to rob, to plunder'.
\end{cousinweb}

\begin{grammarlens}[title={What you've built}]
One more word for telling, explaining and describing at B1.
\end{grammarlens}

\subsection*{Your turn}
\begin{itemize}
\raggedright
  \item \emph{Say it:} \emph{el robo}
  \item \emph{Say it:} \emph{el robo}, once more
  \item \emph{Say it:} a sentence of your own with \emph{el robo}
  \item \emph{Recall:} say the Spanish for the detail, then the Spanish for to steal, then say \emph{el robo} again
\end{itemize}

\begin{tcolorbox}[breakable,colback=teal!4,colframe=teal!35!black,title={Before you move on}]
What does el robo mean? (\textquotedblleft{}The robbery, the theft\textquotedblright{}.) Say it once more.
\end{tcolorbox}
\section[el sospechoso]{el sospechoso — the suspect; suspicious}

\label{lesson:ES-C518-sospechoso}



\begin{quote}
Before the new one: say the Spanish for to steal, then the Spanish for the robbery.
\end{quote}

\subsection*{You'll want to know: el sospechoso}
\textbf{el sospechoso} — \textquotedblleft{}the suspect; suspicious\textquotedblright{}.

\begin{quote}
La policía busca a dos sospechosos.
\end{quote}

\begin{quote}
\emph{The police are looking for two suspects.}
\end{quote}

As an adjective: \emph{un coche sospechoso} — a suspicious car.

\begin{cousinweb}
\emph{sospechoso} is built on \emph{sospecha} 'suspicion' with the ending \emph{-oso}. \emph{Sospecha} is built on \emph{sospechar}, which is inherited from Latin \emph{suspectāre} 'to look up at, to suspect' (EtymDB), from \emph{suspicere}.
\end{cousinweb}

\begin{grammarlens}[title={What you've built}]
One more word for telling, explaining and describing at B1.
\end{grammarlens}

\subsection*{Your turn}
\begin{itemize}
\raggedright
  \item \emph{Say it:} \emph{el sospechoso}
  \item \emph{Say it:} \emph{el sospechoso}, once more
  \item \emph{Say it:} a sentence of your own with \emph{el sospechoso}
  \item \emph{Recall:} say the Spanish for to steal, then the Spanish for the robbery, then say \emph{el sospechoso} again
\end{itemize}

\begin{tcolorbox}[breakable,colback=teal!4,colframe=teal!35!black,title={Before you move on}]
What does el sospechoso mean? (\textquotedblleft{}The suspect; suspicious\textquotedblright{}.) Say it once more.
\end{tcolorbox}
\section[atracar]{atracar — to hold up, to mug; to moor}

\label{lesson:ES-C518-atracar}



\begin{quote}
Before the new one: say the Spanish for the robbery, then the Spanish for the suspect.
\end{quote}

\subsection*{You'll want to know: atracar}
\textbf{atracar} — \textquotedblleft{}to hold up, to mug; to moor\textquotedblright{}.

\begin{quote}
Dos hombres atracaron un banco en el centro.
\end{quote}

\begin{quote}
\emph{Two men held up a bank in the centre.}
\end{quote}

\emph{Un atraco} is a hold-up. For a boat, \emph{atracar} means to moor.

\begin{cousinweb}
The origin of \emph{atracar} is not settled. The sense 'to hold up, to mug' is often linked to the sailors' verb \emph{atracar} 'to bring a boat alongside'. EtymDB traces a Spanish \emph{atracar} to an Arabic verb, without saying which sense it means.
\end{cousinweb}

\begin{grammarlens}[title={What you've built}]
One more word for telling, explaining and describing at B1.
\end{grammarlens}

\subsection*{Your turn}
\begin{itemize}
\raggedright
  \item \emph{Say it:} \emph{atracar}
  \item \emph{Say it:} \emph{atracar}, once more
  \item \emph{Say it:} a sentence of your own with \emph{atracar}
  \item \emph{Recall:} say the Spanish for the robbery, then the Spanish for the suspect, then say \emph{atracar} again
\end{itemize}

\begin{tcolorbox}[breakable,colback=teal!4,colframe=teal!35!black,title={Before you move on}]
What does atracar mean? (\textquotedblleft{}To hold up, to mug; to moor\textquotedblright{}.) Say it once more.
\end{tcolorbox}
\section[huir]{huir — to flee, to run away}

\label{lesson:ES-C518-huir}



\begin{quote}
Before the new one: say the Spanish for the suspect, then the Spanish for to hold up.
\end{quote}

\subsection*{You'll want to know: huir}
\textbf{huir} — \textquotedblleft{}to flee, to run away\textquotedblright{}.

\begin{quote}
El ladrón huyó en un coche rojo.
\end{quote}

\begin{quote}
\emph{The thief fled in a red car.}
\end{quote}

It takes a \emph{y} in many forms: \emph{huyo}, \emph{huyó}, \emph{huyeron}.

\begin{cousinweb}
\emph{huir} is inherited from Latin \emph{fugere} 'to flee' (Lewis \& Short), through old Spanish \emph{fuir} (EtymDB). Latin \emph{f-} became Spanish \emph{h-}. English \emph{fugitive} is from the same verb.
\end{cousinweb}

\begin{grammarlens}[title={What you've built}]
One more word for telling, explaining and describing at B1.
\end{grammarlens}

\subsection*{Your turn}
\begin{itemize}
\raggedright
  \item \emph{Say it:} \emph{huir}
  \item \emph{Say it:} \emph{huir}, once more
  \item \emph{Say it:} a sentence of your own with \emph{huir}
  \item \emph{Recall:} say the Spanish for the suspect, then the Spanish for to hold up, then say \emph{huir} again
\end{itemize}

\begin{tcolorbox}[breakable,colback=teal!4,colframe=teal!35!black,title={Before you move on}]
What does huir mean? (\textquotedblleft{}To flee, to run away\textquotedblright{}.) Say it once more.
\end{tcolorbox}
